%----------------------------------------------------------------------------------------
%	GLOSSARY
%----------------------------------------------------------------------------------------

% Only the used entries will be displayed in the printed list, ie, you need to used a term at least once
% In italic if not in the main document language

% terms definition usage:
% \newglossaryentry{<tag>}{name={<term>},description={<description of the term>}}

\newglossaryentry{gloss}{
name={glossário}, 
description={é uma lista alfabética de termos de um determinado domínio de conhecimento com a definição desses mesmos termos.},
}

\newglossaryentry{pack}{
name={\textit{package}}, 
description={é um ficheiro ou conjunto de ficheiros que contêm comandos \LaTeX{} extra que adicionam novas funcionalidades de estilo ou modificam aquelas já existentes.},
sort={package}	%needed for sorting when using LaTeX commands in the 'name' field
}

\newglossaryentry{lipsum}{
name={\textit{Lorem Ipsum}}, 
description={é uma sequência de palavras, geralmente latinas, utilizada para preencher o espaço destinado a texto numa publicação, por forma a testar as opções de formatação e edição e o arranjo dos elementos gráficos antes da inserção do conteúdo.},
sort={Lorem Ipsum}


% Termos técnicos da dissertação
\newglossaryentry{arbitragem}{
    name={arbitragem energética},
    description={estratégia de operação que consiste em armazenar energia quando o seu custo é baixo e utilizá-la ou vendê-la quando o custo é mais elevado}
}
\newglossaryentry{autoconsumo}{
    name={autoconsumo},
    description={parcela da energia produzida localmente que é consumida na própria instalação}
}
\newglossaryentry{autossuficiencia}{
    name={autossuficiência energética},
    description={parcela do consumo energético da instalação que é satisfeita pela produção local}
}
\newglossaryentry{capacidade_nominal}{
    name={capacidade nominal},
    description={quantidade de carga elétrica que uma bateria pode fornecer nas condições de referência especificadas pelo fabricante}
}
\newglossaryentry{capacidade_util}{
    name={capacidade útil},
    description={quantidade de energia armazenada que pode ser efetivamente utilizada, considerando os limites de operação do sistema}
}
\newglossaryentry{c_rate}{
    name={taxa C},
    description={razão entre a corrente de carga ou descarga e a capacidade nominal da bateria, usada para expressar a rapidez relativa desse processo}
}
\newglossaryentry{ciclo_completo}{
    name={ciclo completo equivalente},
    description={medida de utilização que agrega ciclos parciais de carga e descarga até perfazer a energia correspondente a um ciclo completo}
}
\newglossaryentry{degradacao_calendarial}{
    name={degradação calendarial},
    description={perda irreversível de capacidade ou desempenho que ocorre ao longo do tempo, mesmo quando a bateria não está a ser ciclada}
}
\newglossaryentry{degradacao_ciclica}{
    name={degradação cíclica},
    description={perda irreversível de capacidade ou desempenho associada aos processos de carga e descarga da bateria}
}
\newglossaryentry{eficiencia_ciclo}{
    name={eficiência de ciclo completo},
    description={razão entre a energia fornecida pelo sistema durante a descarga e a energia necessária para o carregar num ciclo completo}
}
\newglossaryentry{estado_carga}{
    name={estado de carga},
    description={estimativa da carga disponível na bateria expressa relativamente à sua capacidade de referência}
}
\newglossaryentry{estado_saude}{
    name={estado de saúde},
    description={indicador do estado de envelhecimento da bateria, frequentemente estimado pela capacidade atual relativamente à capacidade nominal}
}
\newglossaryentry{peak_shaving}{
    name={\textit{peak shaving}},
    description={estratégia que reduz a potência máxima solicitada à rede através da descarga do sistema de armazenamento durante os períodos de maior consumo}
}
\newglossaryentry{sistema_fotovoltaico}{
    name={sistema fotovoltaico},
    description={sistema que converte a energia da radiação solar em energia elétrica por meio de células fotovoltaicas}
}
\newglossaryentry{vida_util_remanescente}{
    name={vida útil remanescente},
    description={tempo ou número de ciclos estimado até o sistema atingir o limite de fim de vida definido para a aplicação}
}
}
